\section{A decision tree with a limit}
\begin{frame}{Decision trees: choices on edges}
\princ{DECISION RULE}
\begin{center}
\begin{tikzpicture}[>=Stealth,every node/.style={draw,rounded corners,fill=blue!8,minimum width=1.55cm,minimum height=0.5cm,font=\small},edge/.style={->,thick}]
\node (s) at (0,2.0) {Start};
\node (e) at (0,0.8) {eligible?};
\node (n) at (-2.2,-0.5) {needs aid?};
\node (x) at (2.2,-0.5) {explain rule};
\node (a) at (-3.1,-1.8) {aid route};
\node (r) at (-1.0,-1.8) {standard route};
\draw[edge] (s)--(e);
\draw[edge] (e)--(n) node[midway,left,font=\scriptsize,draw=none,fill=white] {yes};
\draw[edge] (e)--(x) node[midway,right,font=\scriptsize,draw=none,fill=white] {no};
\draw[edge] (n)--(a) node[midway,left,font=\scriptsize,draw=none,fill=white] {yes};
\draw[edge] (n)--(r) node[midway,right,font=\scriptsize,draw=none,fill=white] {no};
\end{tikzpicture}
\end{center}
\see{yes and no are edge labels; they are not extra vertices.}
\note{Introduce the supplied Tuesday scenario as another concrete tree. Point out every yes/no label on an arrow. The labels describe which branch was taken; the nodes are Start, questions, and outcomes. This is the cleaned decision tree required for the package.}
\end{frame}

\begin{frame}{Trace one route: eligible and needs aid}
\princ{DECISION RULE}
\begin{block}{Input}
eligible = yes; needs aid = yes
\end{block}
\begin{center}\small\texttt{Start -> eligible? -> needs aid? -> aid route}\end{center}
\begin{itemize}
\item Follow the yes edge from eligible? to needs aid?.
\item Follow the yes edge from needs aid? to aid route.
\end{itemize}
\see{the route is a trace through nodes; yes is the label on each selected edge.}
\note{Read the route slowly and ask students to point to the selected arrows. The trace does not include ``yes'' as a node. Make the distinction concrete: the route records visited nodes, while the arrow labels explain each choice.}
\end{frame}

\begin{frame}{The three decision-tree traces}
\princ{TRAVERSAL}
\small
\begin{tabular}{@{}p{3.1cm}p{10.0cm}@{}}
\textbf{preorder} & \texttt{\TraceDecisionPre}\\[5pt]
\textbf{postorder} & \texttt{\TraceDecisionPost}\\[5pt]
\textbf{breadth-first} & \texttt{\TraceDecisionBFS}\\
\end{tabular}
\vspace{8pt}
\mistake{If yes/no appears in a node list, the model has changed: labels became vertices.}
\note{Compare the decision-tree traces with the project-tree traces. The code uses the same traversal functions, but the child order is different: needs aid? comes before explain rule. Ask students to look for yes/no in the lists; they should not find it.}
\end{frame}

\begin{frame}{The language-access omission}
\princ{HIERARCHY LIMIT}
\begin{columns}[T]
\begin{column}{0.48\textwidth}
\begin{block}{What the tree says}
Eligible + needs aid leads to aid route.
\end{block}
\end{column}
\begin{column}{0.48\textwidth}
\begin{alertblock}{What the intake note also says}
\texttt{aid route -> language-access review}
\end{alertblock}
\end{column}
\end{columns}
\vspace{10pt}
\begin{center}\small The tree does not show that relationship.\end{center}
\see{``aid route'' is not proof that language access has been reviewed.}
\note{State the limitation without inventing a new policy. The displayed tree ends at aid route, while the intake note adds a separate relation to language-access review. The tree can be useful and still incomplete. Ask what a stronger conclusion would wrongly assume.}
\end{frame}

\begin{frame}{Bound the claim}
\princ{CHECK THE WORK}
\begin{block}{Supported claim}
For this displayed tree, this child order, and this input, the trace reaches aid route.
\end{block}
\begin{alertblock}{Unsupported stronger claim}
Everyone routed to aid can use the aid route immediately.
\end{alertblock}
\begin{itemize}
\item The missing relation may need another node: language-access review.
\item Or it must remain an explicit limitation of the model.
\end{itemize}
\note{Have students compare the bounded and unbounded claims. The first is directly checked by a trace. The second needs information not represented in the tree. Do not rewrite the Decision Gate or turn the omission into a real-person example; keep it synthetic and model-focused.}
\end{frame}
